\chapter{Theory}
\label{chap:theory}

% ------------------------------------------------------------------
% STRUCTURE (restructured 2026-09-30, inspired by Nordal (2026), Ch. 2)
%   2.1 Planar Snake Robot Model
%   2.2 Grasping with Friction
%       2.2.1 Grasp Matrix
%       2.2.2 Hand Jacobian Matrix
%       2.2.3 Contact and Friction Model
%       2.2.4 Coulomb Friction Cone
%   2.3 Force Closure
%       2.3.1 Nullspace Cases for the Grasp Matrix and Hand Jacobian
%       2.3.2 Frictional Form Closure  (LP2)
%       2.3.3 Force Closure            (LP3)
%
% Provenance: every moved block is marked with a comment
%   % [Moved from <file>, <section>]
% Text that was removed because it duplicated something kept here is
% collected, with its origin, in chapters/removed_duplicates.tex
% (not \input in main.tex).
%
% Original chapter TODO list (kept for reference):
%  - Kinematics/dynamics model of the snake robot relevant to
%    contact forces (link model, contact points, friction model
%    e.g. Coulomb friction cones).
%  - Contact wrench space, wrench space formulation for a set of
%    contacts between the robot and an object/obstacle.
%  - Formal definition of force closure (and form closure) for a
%    set of contacts, and the standard convex-hull/wrench-space
%    tests for it.
%  - Any coordinate frames, notation, and conventions used
%    consistently for the rest of the report -- define them here
%    once so later chapters can just reference \cref{chap:theory}.
%    [The empty section "Notation and Conventions" was removed; put
%     notation in a List of Symbols in the front matter instead, as
%     Nordal (2026) does.]
% ------------------------------------------------------------------

% ==================================================================
\section{Planar Snake Robot Model}
\label{sec:snake-model}
% ==================================================================

% [Moved from margin_force_closure.tex, section "Properties of the
%  snake robot"]
Should follow the definition in \textcite{gravdahl2026formclosure}.
As in \cite{gravdahl2026formclosure,nordal2026grip} the link width is
neglected; for frictionless contacts this is exact, while with friction it
slightly changes the moment arms of the tangential contact forces.

\todo[inline]{Write the planar snake model. Cf.\ Nordal (2026) Sec.~2.1
(unconstrained and constrained kinematics): $n$ links of length $l$, joint
angles $\boldsymbol{\phi}\in\mathbb{R}^{n-1}$, base coordinates
$\boldsymbol{q}_0=[x_0\ y_0\ \phi_0]^\top$ and configuration vector
$\boldsymbol{q}\in\mathbb{R}^{n+2}$.}

% ==================================================================
\section{Grasping with Friction}
\label{sec:grasping}
% ==================================================================

\subsection{Grasp Matrix}
\label{sec:grasp-matrix}

% [Moved from background.tex, section "The Grasp Matrix", subsection
%  "Definition following Gravdahl et al. (2026)"]
\subsubsection{Definition Following Gravdahl et al.}
\label{sec:grasp-matrix-gravdahl}

\textcite{gravdahl2026formclosure} interpret a planar snake robot in contact
with a set of obstacles as a grasp, where the obstacles play the role of the
``fingers'' and the snake's body is the grasped object. The grasp matrix
$\boldsymbol{G}$ relates the object twist and the contact-point twists,
\begin{equation}
    \boldsymbol{\nu}_{\text{contacts}} = \boldsymbol{G}^\top \boldsymbol{\nu}_{\text{object}} ,
    \label{eq:grasp-matrix-velocity}
\end{equation}
and, dually, maps the contact forces to the resultant wrench acting on the
object,
\begin{equation}
    \boldsymbol{w}_{\text{object}} = \boldsymbol{G}\, \boldsymbol{f}_{\text{contacts}} .
    \label{eq:grasp-matrix-wrench}
\end{equation}

For a single contact $i$, the \emph{partial grasp matrix} is
\begin{equation}
    \tilde{\boldsymbol{G}}_i^\top = \bar{\boldsymbol{R}}_i^\top \boldsymbol{P}_i^\top \in \mathbb{R}^{3 \times 3} ,
    \label{eq:partial-grasp-matrix}
\end{equation}
where
\begin{equation}
    \bar{\boldsymbol{R}}_i =
    \begin{bmatrix}
        \boldsymbol{R} & \boldsymbol{0} \\
        \boldsymbol{0} & 1
    \end{bmatrix} \in \mathbb{R}^{3\times3},
    \qquad
    \boldsymbol{P}_i =
    \begin{bmatrix}
        \boldsymbol{I}_{2\times2} & \boldsymbol{0} \\
        \boldsymbol{S}(\boldsymbol{p}_i - \boldsymbol{p}_0) & 1
    \end{bmatrix} \in \mathbb{R}^{3\times3} ,
    \label{eq:Rbar-P}
\end{equation}
with $\boldsymbol{R}$ the $2\times2$ rotation matrix, $\boldsymbol{I}_{2\times2}$ the $2\times2$
identity matrix, and $\boldsymbol{S}(\boldsymbol{p}_i - \boldsymbol{p}_0)$ the skew-symmetric
representation of the displacement between contact point $i$ and a chosen
reference point $\boldsymbol{p}_0$,
\begin{equation}
    \boldsymbol{S}(\boldsymbol{a}) = \begin{bmatrix} -a_y & a_x \end{bmatrix}, \qquad \boldsymbol{a} = (a_x, a_y)^\top .
    \label{eq:skew-2d}
\end{equation}

Stacking the partial grasp matrices of all $n_c$ contacts gives the
\emph{complete} grasp matrix, which is then reduced by the contact model
matrix $\boldsymbol{H}$:
\begin{equation}
    \tilde{\boldsymbol{G}}^\top = \begin{bmatrix} \tilde{\boldsymbol{G}}_1 & \cdots & \tilde{\boldsymbol{G}}_{n_c} \end{bmatrix}^\top \in \mathbb{R}^{3n_c \times 3} ,
    \qquad
    \boldsymbol{G}^\top = \boldsymbol{H}\,\tilde{\boldsymbol{G}}^\top .
    \label{eq:complete-grasp-matrix}
\end{equation}
Gravdahl et al.\ restrict their analysis to the point-contact-without-friction
(PwoF) model, for which $\boldsymbol{H}$ keeps only the normal-force component
of each contact:
\begin{equation}
    \boldsymbol{H}^{\mathrm{PwoF}} =
    \begin{bmatrix}
        1 & 0 & 0 & 0 & 0 & 0 & \cdots & 0 & 0 & 0 \\
        0 & 0 & 0 & 1 & 0 & 0 & \cdots & 0 & 0 & 0 \\
        \vdots & & & & & & \ddots & & & \vdots \\
        0 & 0 & 0 & 0 & 0 & 0 & \cdots & 1 & 0 & 0
    \end{bmatrix} \in \mathbb{R}^{n_c \times 3n_c} ,
    \label{eq:H-pwof-gravdahl}
\end{equation}
so that $\boldsymbol{G} \in \mathbb{R}^{3 \times n_c}$: one column per contact,
along the contact normal.

% [Moved from background.tex, section "The Grasp Matrix", subsection
%  "Generalization: the Handbook of Robotics framework"]
\subsubsection{Generalization to Frictional Contacts}
\label{sec:grasp-matrix-handbook}

Equations~\eqref{eq:grasp-matrix-velocity}--\eqref{eq:complete-grasp-matrix} are, in
fact, a specialization of the general grasp-analysis framework in the
\emph{Handbook of Robotics}, Ch.~28 \cite{prattichizzo2008grasping}
(Gravdahl et al.\ adopt this nomenclature directly). The Handbook treats $\boldsymbol{H}$
as a free choice determined by the \emph{contact model}
(\cref{sec:contact-model}), so that form closure
(frictionless, PwoF) and force closure (frictional, HF) fall out of the same
construction --- only $\boldsymbol{H}$ changes. This is the natural route from
Gravdahl et al.'s form-closure analysis to the force-closure margin this
thesis develops: the geometry ($\tilde{\boldsymbol{G}}_i$, $\bar{\boldsymbol{R}}_i$,
$\boldsymbol{P}_i$) is untouched; only the contact-model selection matrix
$\boldsymbol{H}_i$ needs to be replaced.

For a single planar contact, the Handbook (Table~28.3) gives
\begin{equation}
    \boldsymbol{H}_i^{\mathrm{PwoF}} = \begin{bmatrix} 1 & 0 & 0 \end{bmatrix} \in \mathbb{R}^{1\times3} , \qquad \ell_i = 1 ,
    \label{eq:Hi-pwof-planar}
\end{equation}
\begin{equation}
    \boldsymbol{H}_i^{\mathrm{HF}} = \begin{bmatrix} \boldsymbol{I}_{2\times2} & \boldsymbol{0}_{2\times1} \end{bmatrix} \in \mathbb{R}^{2\times3} , \qquad \ell_i = 2 ,
    \label{eq:Hi-hf-planar}
\end{equation}
noting that in the planar case the hard-finger (HF) and soft-finger (SF)
models coincide, since there is no out-of-plane axis for a torsional moment
to act about. Equation~\eqref{eq:Hi-pwof-planar} is exactly
Gravdahl et al.'s single-contact row in~\eqref{eq:H-pwof-gravdahl}.

\paragraph{Adapting to $n_c$ contacts.}
Stacking~\eqref{eq:Hi-pwof-planar} and~\eqref{eq:Hi-hf-planar} block-diagonally
over all $n_c$ contacts gives the complete contact model matrix for each
case:
\begin{equation}
    \boldsymbol{H}^{\mathrm{PwoF}} = \mathrm{Blockdiag}\big(\boldsymbol{H}_1^{\mathrm{PwoF}}, \ldots, \boldsymbol{H}_{n_c}^{\mathrm{PwoF}}\big) \in \mathbb{R}^{n_c \times 3n_c} ,
    \label{eq:H-pwof-nc}
\end{equation}
\begin{equation}
    \boldsymbol{H}^{\mathrm{HF}} = \mathrm{Blockdiag}\big(\boldsymbol{H}_1^{\mathrm{HF}}, \ldots, \boldsymbol{H}_{n_c}^{\mathrm{HF}}\big) \in \mathbb{R}^{2n_c \times 3n_c} ,
    \label{eq:H-hf-nc}
\end{equation}
written out explicitly, the hard-finger case is
\begin{equation}
    \boldsymbol{H}^{\mathrm{HF}} =
    \left[\begin{array}{ccc|ccc|c|ccc}
        1 & 0 & 0 & 0 & 0 & 0 & \cdots & 0 & 0 & 0 \\
        0 & 1 & 0 & 0 & 0 & 0 & \cdots & 0 & 0 & 0 \\
        0 & 0 & 0 & 1 & 0 & 0 & \cdots & 0 & 0 & 0 \\
        0 & 0 & 0 & 0 & 1 & 0 & \cdots & 0 & 0 & 0 \\
        \vdots & & & & & & \ddots & & & \vdots \\
        0 & 0 & 0 & 0 & 0 & 0 & \cdots & 1 & 0 & 0 \\
        0 & 0 & 0 & 0 & 0 & 0 & \cdots & 0 & 1 & 0
    \end{array}\right] \in \mathbb{R}^{2n_c \times 3n_c} .
    \label{eq:H-hf-nc-explicit}
\end{equation}

Applying $\boldsymbol{H}$ to the complete partial grasp matrix from
\eqref{eq:complete-grasp-matrix} (unchanged from the Gravdahl derivation)
gives the final grasp matrix for each contact model:
\begin{align}
    \boldsymbol{G}^{\mathrm{PwoF}} &\in \mathbb{R}^{3 \times n_c} , &
    \boldsymbol{G}^{\mathrm{PwoF}\top} &= \boldsymbol{H}^{\mathrm{PwoF}}\,\tilde{\boldsymbol{G}}^\top ,
    \label{eq:G-pwof-nc}\\
    \boldsymbol{G}^{\mathrm{HF}} &\in \mathbb{R}^{3 \times 2n_c} , &
    \boldsymbol{G}^{\mathrm{HF}\top} &= \boldsymbol{H}^{\mathrm{HF}}\,\tilde{\boldsymbol{G}}^\top .
    \label{eq:G-hf-nc}
\end{align}

Under the hard-finger model, each contact contributes \emph{two} columns to
$\boldsymbol{G}$ instead of one: the normal-force direction (same as PwoF)
and the tangential (friction) direction. This is what allows the contacts to
generate wrenches off the contact normals, which is a necessary ingredient
for force closure, but is absent by construction from
Gravdahl et al.'s PwoF-only form-closure analysis. The two force
components at a hard-finger contact are not independent: they are coupled by
the Coulomb friction cone $|f_{\mathrm{t},i}| \le \mu_i f_{\mathrm{n},i}$
(\cref{sec:friction-cone}), so
for the force-closure margin computation in
Chapter~\ref{chap:margin} each contact's two columns are replaced by the two
edge generators of the (linearized) friction cone before imposing
$\boldsymbol{\lambda} \ge \boldsymbol{0}$.

% [Moved from background.tex, section "Wrenches, Contact Forces, and
%  Internal Forces", subsections "Contact Forces and Wrenches" and
%  "Wrench Acting on the Object" (merged). Changed: wrench written for
%  the planar case (R^3 instead of R^6), and the broken LaTeX in the
%  lambda vector fixed. The repeated equations w = G lambda were
%  removed -> removed_duplicates.tex.]
\subsubsection{Contact Forces and Wrenches}
\label{sec:contact-forces-wrenches}

To analyze force closure, it is necessary to distinguish between the forces
applied at the contacts, the resulting wrench acting on the object, and
internal forces within the contact force system. These concepts are
connected through the grasp matrix.

Let the contact forces transmitted at the $n_c$ contact points be stacked in
the vector
\begin{equation}
    \boldsymbol{\lambda}
    =
    \begin{bmatrix}
        \boldsymbol{\lambda}_1^\top & \boldsymbol{\lambda}_2^\top & \cdots & \boldsymbol{\lambda}_{n_c}^\top
    \end{bmatrix}^\top ,
    \label{eq:lambda-stacked}
\end{equation}
which for HF contacts in the plane is ordered as
$\boldsymbol{\lambda} = [f_{n1},\, f_{t1},\, f_{n2},\, f_{t2},\, \ldots,\, f_{n n_c},\, f_{t n_c}]^\top$,
where $f_{ni}$ and $f_{ti}$ denote the normal and tangential force
components at contact~$i$.
% [The explicit ordering sentence above is moved from implementation.tex,
%  section "Algorithm".]
The grasp matrix maps the contact forces to the resulting wrench acting on
the object,
\begin{equation}
    \boldsymbol{w} = \boldsymbol{G}\boldsymbol{\lambda},
    \qquad
    \boldsymbol{w} =
    \begin{bmatrix}
        \boldsymbol{f} \\ \tau
    \end{bmatrix}
    \in \mathbb{R}^{3},
    \label{eq:wrench-map}
\end{equation}
where $\boldsymbol{f}=(f_x,f_y)$ is the net force acting on the object and
$\tau$ the net torque about the reference point $\boldsymbol{p}_0$.
Each column of $\boldsymbol{G}$ is therefore a \emph{primitive wrench} ---
the wrench produced by a unit force along one basis direction (normal or
tangential) at one contact.
% [Previous two sentences moved from margin_force_closure.tex,
%  subsection "Wrench Space and Force Closure".]
Consequently, the grasp matrix describes how the individual contact forces
contribute to the overall force and torque acting on the object.

For example, consider two contact forces acting on opposite sides of an
object. If the forces act in the same direction, they can produce a net
translational force. If the forces are arranged such that their
translational effects cancel while their moment contributions do not, they
can instead produce a torque.

% ------------------------------------------------------------------
\subsection{Hand Jacobian Matrix}
\label{sec:hand-jacobian}

\todo[inline]{Write this subsection (new -- no existing text). Cf.\ Nordal
(2026) Sec.~2.4.4--2.4.5 (task/contact Jacobians and the hand/snake
Jacobian) and Handbook Ch.~28.2. Key points from the project notes:
kinematics $\boldsymbol{J}\dot{\boldsymbol{q}}=\boldsymbol{G}^\top\boldsymbol{\nu}$,
statics $\boldsymbol{\tau}=\boldsymbol{J}^\top\boldsymbol{\lambda}$; for the
floating-base snake $\boldsymbol{J}_{\mathrm{snk}}=[\boldsymbol{G}^\top\ |\ \boldsymbol{J}_\phi]$.}

% ------------------------------------------------------------------
% [Moved from theory.tex, section "Contact and Friction Model".
%  NOTE: the HF/SF paragraphs are close to verbatim from the Handbook
%  (pp. 675-676) -- paraphrase before submission. The Handbook
%  paragraph on internal forces (p. 678) that was here has been moved
%  to Sec. "Nullspace Cases" below.]
\subsection{Contact and Friction Model}
\label{sec:contact-model}

A hard-finger (HF) model is used when there is significant contact friction,
but the contact patch is small, so that no appreciable friction moment
exists. When this model is applied to a contact, all three translational
velocity components of the contact point on the hand (i.e., the first three
components of $\boldsymbol{\nu}_{i,\mathrm{hnd}}$) and all three components of
the contact force are transmitted through the contact. None of the angular
velocity components or moment components are transmitted.
\textcolor{red}{Claude: For a snake robot pressing a body segment against an
obstacle, HF is the used default (point/line contact, negligible torque about
the surface normal).}

The soft-finger (SF) model is used in situations in which the surface
friction and the contact patch are large enough to generate significant
friction forces and a friction moment about the contact normal. At a contact
where this model is enforced, the three translational velocity components of
the contact on the hand and the angular velocity component about the contact
normal are transmitted (i.e., the first four components of
$\boldsymbol{\nu}_{i,\mathrm{hnd}}$). Similarly, all three components of
contact force and the normal component of the contact moment are transmitted
\cite[p.~675]{prattichizzo2008grasping}.

In planar cases, the SF and HF models are equivalent, because the object and
the hand lie in a plane. Rotations about the contact normals would cause
out-of-plane motions \cite[p.~676]{prattichizzo2008grasping}.

% ------------------------------------------------------------------
% [Moved from theory.tex, section "Coulomb Friction Cone"]
\subsection{Coulomb Friction Cone}
\label{sec:friction-cone}

The Coulomb friction cone is defined as in \cite{prattichizzo2008grasping}.
At a contact point $i$, the contact force is decomposed into a normal
component $f_{in}$ (pressing into the surface) and two tangential (friction)
components $f_{it}, f_{io}$ lying in the contact plane. Coulomb's friction
law states that the tangential force cannot exceed the normal force scaled by
the friction coefficient $\mu_i$, without the contact slipping:
\begin{equation}
    \mathcal{F}_i = \left\{ (f_{in}, f_{it}, f_{io}) \;\middle|\; \sqrt{f_{it}^2 + f_{io}^2} \le \mu_i f_{in} \right\}.
    \label{eq:friction-cone-explicit}
\end{equation}
The set $\mathcal{F}_i$ contains all admissible (non-slipping) contact forces
at point $i$. Plotted with $f_{in}$ as one axis and $(f_{it}, f_{io})$ as the
other two, this set forms a \emph{cone}: the allowed friction magnitude grows
linearly with the normal force.

More generally, when all $\ell_i$ force components at contact $i$ are stacked
into a vector $\boldsymbol{\lambda}_i \in \mathbb{R}^{\ell_i}$, the friction
cone is written compactly using a (possibly weighted) norm $\|\cdot\|_w$:
\begin{equation}
    \mathcal{F}_i = \left\{ \boldsymbol{\lambda}_i \in \mathbb{R}^{\ell_i} \;\middle|\; \|\boldsymbol{\lambda}_i\|_w \le \mu_i f_{in} \right\},
    \label{eq:friction-cone-general}
\end{equation}
where $\|\boldsymbol{\lambda}_i\|_w$ denotes a weighted quadratic norm;
\cref{tab:norms-contact-models} defines useful norms for the different
contact models.
\begin{table}[htbp]
    \centering
    \caption{Norms for the three main contact models.}
    \label{tab:norms-contact-models}
    \begin{tabular}{ll}
        \toprule
        Model & $\|\boldsymbol{\lambda}_i\|_w$ \\
        \midrule
        PwoF & $0$ \\[4pt]
        HF   & $\dfrac{1}{\mu_i}\sqrt{f_{it}^2 + f_{io}^2}$ \\[10pt]
        SF   & $\dfrac{1}{\mu_i}\sqrt{f_{it}^2 + f_{io}^2 + a v_i^{-1} \lvert m_{in} \rvert^2}$ \\
        \bottomrule
    \end{tabular}
\end{table}

% [Moved from margin_force_closure.tex, subsection "Wrench Space and
%  Force Closure" (planar contact description and eq. "friction-cone").
%  Symbols changed from mu to mu_i to match the general cone above.]
\paragraph{Planar case.}
For a planar contact between the snake robot's body and an obstacle at
$\boldsymbol{p}_i \in \mathbb{R}^2$, with unit outward normal
$\boldsymbol{n}_i$ and tangent $\boldsymbol{t}_i = \boldsymbol{n}_i^{\perp}$,
the out-of-plane component $f_{io}$ vanishes and the cone reduces to a planar
wedge. The contact transmits the force
\begin{equation}
    \boldsymbol{f}_i = f_{n,i}\,\boldsymbol{n}_i + f_{t,i}\,\boldsymbol{t}_i,
    \qquad f_{n,i} \ge 0,\ \ |f_{t,i}| \le \mu_i f_{n,i},
    \label{eq:friction-cone}
\end{equation}
but no moment about the contact point (HF contact model).

\textbf{Interpretation:} $\mathcal{F}_i$ is the set of physically admissible
contact forces that keep contact $i$ from slipping. A required contact force
lying \emph{inside} the cone means the contact holds; one falling outside
means the surfaces slip. Force closure requires a set of contacts whose
friction cones, in combination, can resist any external disturbance
force/wrench \cite[p.~685]{prattichizzo2008grasping}.

% [Moved from theory.tex, section "Force Closure Tests" (intro and all
%  its subsections). Placed here because the linearized cone (matrix F)
%  is what LP2 and LP3 below are built on. Subsections demoted to
%  paragraphs.]
\subsubsection{Linearized Friction Cone}
\label{sec:linearized-cone}

The friction cone is a smooth, curved surface: the friction constraint
\eqref{eq:friction-cone-explicit} is quadratic and requires nonlinear
programming techniques to solve. It is, however, possible to work with
approximate tests \cite[p.~686]{prattichizzo2008grasping}. One approach is to
replace it with a polyhedral (pyramidal) cone, defined as the nonnegative
span of a finite set of generator vectors. This way it is possible to solve
it with linear programming techniques. This will later in the report be used
for checking whether a candidate contact configuration for the snake robot
satisfies force closure.

\paragraph{Generators of a polyhedral cone.}
Any polyhedral cone can be represented as the set of all nonnegative linear
combinations of a finite number $n_g$ of generator vectors
$\mathbf{s}_{i1}, \dots, \mathbf{s}_{i n_g}$. A force $\mathbf{f}$ lying
inside the approximate cone at contact $i$ can be written explicitly as
\begin{equation}
    \mathbf{f} = \sum_{k=1}^{n_g} \sigma_k \, \mathbf{s}_{ik}, \qquad \sigma_k \ge 0 ,
    \label{eq:generator-combination}
\end{equation}
where each $\sigma_k$ is a nonnegative generator weight. Geometrically, the
$n_g$ generators are the edges of a pyramid inscribed in (or circumscribing)
the true circular friction cone; increasing $n_g$ improves the
approximation.

Equation~\eqref{eq:generator-combination} can be written compactly for the
full set of contact forces $\boldsymbol{\lambda}_i$ transmitted at contact
$i$ as
\begin{equation}
    G_i \boldsymbol{\lambda}_i = S_i \boldsymbol{\sigma}_i, \qquad \boldsymbol{\sigma}_i \ge 0 ,
    \label{eq:generator-matrix-form}
\end{equation}
where $S_i = (\mathbf{s}_{i1} \; \cdots \; \mathbf{s}_{i n_g})$ collects the
generators as columns, and $\boldsymbol{\sigma}_i \ge 0$ is a vector of
nonnegative weights. For a frictionless point contact (PwoF), $n_g = 1$ and
$S_i$ reduces to the single contact normal direction.
% TODO: in the planar HF case the wedge (eq:friction-cone) is exactly
% represented by n_g = 2 generators, n_i +/- mu_i t_i -- state this.

\paragraph{Generators for the hard-finger (HF) model.}
For a hard-finger contact, the friction cone is approximated by the
nonnegative span of $n_g$ uniformly spaced contact-force generators, giving
an inscribed regular polyhedral cone (\cref{fig:polyhedral-cone}). The
generator matrix is
\begin{equation}
    S_i =
    \begin{pmatrix}
        \cdots & 1 & \cdots \\
        \cdots & \mu_i \cos\!\left(\dfrac{2k\pi}{n_g}\right) & \cdots \\
        \cdots & \mu_i \sin\!\left(\dfrac{2k\pi}{n_g}\right) & \cdots
    \end{pmatrix}, \qquad k = 1, \dots, n_g .
    \label{eq:Si-HF}
\end{equation}
The top row fixes a unit normal component for every generator, while the
remaining two rows place the $n_g$ generators evenly around a circle of
radius $\mu_i$ in the tangential $(f_{it}, f_{io})$-plane. If a
\emph{circumscribing} polyhedral cone is preferred instead of an inscribed
one, $\mu_i$ in \eqref{eq:Si-HF} is simply replaced by
$\mu_i / \cos(\pi/n_g)$.

\begin{figure}[htbp]
    \centering
    % Replace with your own figure/drawing of the linearized (polyhedral) friction cone,
    % analogous to Fig. 28.10 in the Springer Handbook of Robotics.
    \includegraphics[width=0.55\linewidth]{Figurer/ntnulogo.pdf}
    \caption{Quadratic (Coulomb) friction cone approximated as a polyhedral cone with seven generators.}
    \label{fig:polyhedral-cone}
\end{figure}

\paragraph{Generators for the soft-finger (SF) model.}
For a soft-finger contact, the torsional friction component is decoupled
from the tangential components, so each generator is extended to four
dimensions:
\begin{equation}
    \mathbf{s}_{ik} = \begin{pmatrix} 1 & 0 & 0 & \pm b \upsilon_i \end{pmatrix}^\top ,
\end{equation}
where $b$ is a characteristic length used to keep units consistent between
force and torsional-moment components. Collecting these generators gives
\begin{equation}
    S_i =
    \begin{pmatrix}
        \cdots & 1 & \cdots & 1 & 1 \\
        \cdots & \mu_i \cos(2k\pi/n_g) & \cdots & 0 & 0 \\
        \cdots & \mu_i \sin(2k\pi/n_g) & \cdots & 0 & 0 \\
        \cdots & 0 & \cdots & b\upsilon_i & -b\upsilon_i
    \end{pmatrix}.
    \label{eq:Si-SF}
\end{equation}
% NOTE: SF = HF in the plane (Sec. Contact and Friction Model), so this
% paragraph can probably be cut or shortened to one sentence.

\paragraph{Assembling the full grasp.}
The set of all contact wrenches that the hand can apply without violating
the friction law at any contact is written as
\begin{equation}
    G \boldsymbol{\lambda} = S \boldsymbol{\sigma}, \qquad \boldsymbol{\sigma} \ge 0 ,
\end{equation}
where $S = (S_1, \dots, S_{n_c})$ stacks the per-contact generator matrices
and $\boldsymbol{\sigma} = (\boldsymbol{\sigma}_1^\top, \dots,
\boldsymbol{\sigma}_{n_c}^\top)^\top$ collects all generator weights, with
$n_c$ the number of contacts.

\paragraph{Dual (face-normal) form.}
It is often convenient to reformulate the friction constraint in dual form,
using the \emph{face normals} of the polyhedral cone rather than its
generators:
\begin{equation}
    F_i \boldsymbol{\lambda}_i \ge 0 ,
    \label{eq:face-normal-form}
\end{equation}
where each row of $F_i$ is normal to a face formed by two adjacent generators
of the approximate cone. For an HF contact, row $k$ of $F_i$ can be computed
directly as the cross product of consecutive generators,
$\mathbf{s}_{ik} \times \mathbf{s}_{i,k+1}$. For an SF contact the generators
are four-dimensional, so simple cross products no longer suffice, and a
general generator-to-face-normal conversion method must be used instead.

Combining the per-contact constraints \eqref{eq:face-normal-form} over all
$n_c$ contacts gives the compact form
\begin{equation}
    F \boldsymbol{\lambda} \ge 0 , \qquad F = \operatorname{Blockdiag}(F_1, \dots, F_{n_c}) .
    \label{eq:face-normal-full}
\end{equation}

% ==================================================================
\section{Force Closure}
\label{sec:force-closure}
% ==================================================================

% [Moved from theory.tex, section "Force Closure: Formal Definition"
%  (first paragraph). Informal note -- rewrite as prose.]
\textcolor{red}{Claude: Most sources use ``force closure'' for a contact-only
property, which is exactly the Handbook's frictional form closure
(Def.~28.9). Only a smaller group, including the Handbook's own Def.~28.10,
makes force closure depend on the robot's actuators. The Handbook itself says
Def.~28.9 is what Murray et al.\ call force closure.}

% [Moved from margin_force_closure.tex, subsection "Wrench Space and
%  Force Closure" (definition of the wrench space). Contact force
%  vector renamed c -> lambda to match the rest of the chapter.]
The set of all wrenches the contact configuration can actually produce, once
the friction cone constraints~\eqref{eq:friction-cone} are imposed on
$\boldsymbol{\lambda}$, is the \emph{wrench space} (or wrench cone) of the
grasp:
\begin{equation}
    \mathcal{W} = \left\{\, \boldsymbol{G}\boldsymbol{\lambda} \ \middle|\ \boldsymbol{\lambda}_i \in \mathcal{F}_i \text{ for all } i \,\right\} \subset \mathbb{R}^3.
    \label{eq:wrench-space}
\end{equation}
Because the friction-cone constraints are conic (positively homogeneous),
$\mathcal{W}$ is a convex cone spanned by the primitive wrenches, obtained as
the Minkowski sum of the (linearized) friction cones mapped through
$\boldsymbol{G}$.

% [Moved from background.tex, subsection "Relation to the Force-Closure
%  Test". Stray markdown ``` fences removed (they broke compilation);
%  the boxed lambda -> w summaries at the end were duplicates and are in
%  removed_duplicates.tex.]
The approximate force-closure test used in this work consists of three
stages \cite[Sec.~28.4.2]{prattichizzo2008grasping}:
\begin{enumerate}
    \item \textbf{Rank of the grasp matrix:}
    \begin{equation}
        \operatorname{rank}(\boldsymbol{G})=n_{\nu}.
    \end{equation}
    This determines whether the contacts can generate all required wrench
    directions (\cref{sec:nullspace-cases}).

    \item \textbf{Frictional form closure:}
    LP2 determines whether there exists a contact-force configuration that
    produces no net wrench while remaining strictly inside the friction
    constraints. A positive optimal value, $d^*>0$, indicates frictional form
    closure (\cref{sec:lp2}).

    \item \textbf{Control of internal forces:}
    LP3 examines whether an internal force exists that is also uncontrollable
    according to the robot's Jacobian. If $d^*>0$, such an internal force
    exists and force closure does not hold. If $d^*=0$, the required
    force-closure condition is satisfied, assuming the previous conditions
    have been fulfilled (\cref{sec:lp3}).
\end{enumerate}

% ------------------------------------------------------------------
\subsection{Nullspace Cases for the Grasp Matrix and Hand Jacobian}
\label{sec:nullspace-cases}

% [Moved from background.tex, subsection "Wrench Directions". Changed:
%  planar case (n_nu = 3) stated explicitly.]
\paragraph{Rank of $\boldsymbol{G}$ and wrench directions.}
A wrench direction represents a possible direction in the object's wrench
space. In three-dimensional space, a wrench has six independent components:
three translational force components and three rotational torque
components; in the plane it has three, $(f_x, f_y, \tau)$. Examples of
wrench directions include:
\begin{itemize}
    \item a force in the positive $x$-direction,
    \item a force in the negative $y$-direction,
    \item a torque about the $z$-axis, or
    \item a combination of forces and torques.
\end{itemize}
Thus, the wrench space can be viewed as the space containing all possible
combinations of forces and torques acting on the object.

For a grasp to have the ability to resist arbitrary external disturbances,
the contacts must be capable of generating wrenches in all relevant
directions. This is examined through the rank of the grasp matrix,
\begin{equation}
    \operatorname{rank}(\boldsymbol{G}) = n_{\nu},
    \label{eq:rank-condition}
\end{equation}
where $n_{\nu}$ denotes the dimension of the object's wrench space
($n_\nu = 6$ in three dimensions, $n_\nu = 3$ in the plane). If the rank is
smaller than $n_{\nu}$, there exists at least one wrench direction that
cannot be generated by the contact forces. Consequently, the contacts cannot
provide force closure. The rank condition can therefore be interpreted as:
\emph{can the contact forces generate wrenches in every required direction
of the object's wrench space?}

% [Moved from background.tex, subsection "Internal Forces"; merged with
%  the Handbook p. 678 paragraph on internal forces from theory.tex,
%  section "Contact and Friction Model".]
\paragraph{Internal forces, $\mathcal{N}(\boldsymbol{G})$.}
Not all contact forces necessarily result in a net wrench on the object.
Some combinations of contact forces cancel each other such that
\begin{equation}
    \boldsymbol{G}\boldsymbol{\lambda} = \mathbf{0}.
    \label{eq:internal-force}
\end{equation}
Such force combinations are referred to as \textit{internal forces}. They
produce no net force or torque on the object, but they can still generate
significant forces at the individual contacts. For example, consider a robot
squeezing an object from two opposite sides, $\mathbf{F}_1 = -\mathbf{F}_2$.
The two forces cancel with respect to the object, $\mathbf{F}_1+\mathbf{F}_2
= \mathbf{0}$, and assuming that their moments also cancel, the
corresponding contact-force vector satisfies~\eqref{eq:internal-force}. The
object therefore experiences no net wrench, even though the robot is applying
forces at the contacts.

The set of all contact-force vectors that produce no net wrench is the
nullspace of the grasp matrix:
\begin{equation}
    \mathcal{N}(\boldsymbol{G})
    = \left\{
    \boldsymbol{\lambda}
    \mid
    \boldsymbol{G}\boldsymbol{\lambda}=\mathbf{0}
    \right\}.
\end{equation}
These wrench intensities do not contribute to the acceleration of the
object; instead, they affect the tightness of the grasp. Internal forces
therefore play a fundamental role in maintaining grasps that rely on
friction \cite[p.~678]{prattichizzo2008grasping}: they can increase the
normal forces at the contacts and thereby increase the available frictional
forces without directly accelerating the object.

% [Moved from theory.tex, section "Nullspace cases for G and J"
%  (notes -- expand into prose).]
\paragraph{Grasp classes.}
From the Handbook of Robotics \cite{prattichizzo2008grasping}:
\textbf{graspable, indeterminate, defective, redundant, hyperstatic}.
\textcolor{red}{``Because the snake's base is unactuated and unsupported, its
Jacobian contains $\boldsymbol{G}$ as the base block, and the hand-grasp
distinction between defective and hyperstatic grasps collapses: a contact
force that loads no joint and no base coordinate is necessarily internal.''}
% TODO: one short paragraph per class, with the nullspace that defines it:
%   graspable N(G), indeterminate N(G^T), redundant N(J),
%   defective N(J^T), hyperstatic N(G) cap N(J^T).

% ------------------------------------------------------------------
\subsection{Frictional Form Closure}
\label{sec:frictional-form-closure}

% [Moved from theory.tex, section "Force Closure: Formal Definition"
%  (second paragraph, first half).]
Def.~28.9 in \cite{prattichizzo2008grasping} requires
$\operatorname{rank}\boldsymbol{G} = n_\nu$ and a squeeze
$\boldsymbol{\lambda} \in \operatorname{Int}(\mathcal{F})$ with
$\boldsymbol{G}\boldsymbol{\lambda} = \mathbf{0}$ (contacts and friction
only).

\subsubsection{LP2}
\label{sec:lp2}

% [Merged from two places:
%  - margin_force_closure.tex, subsection "LP2" (first sentences);
%  - implementation.tex, section "Algorithm" (LP formulation and the
%    explanation of d and of the normalization).
%  The implementation.tex version of the LP is kept because its sign
%  convention (F lambda >= 1 d) matches the face-normal form
%  eq:face-normal-full. The margin-chapter version (F c <= -1 d, with
%  variable c) is in removed_duplicates.tex.]
Following the frictional form-closure test of
\cite[Sec.~28.4.2]{prattichizzo2008grasping}, the friction cone at each
contact is linearized into the half-space constraints
$F\boldsymbol{\lambda} \ge \mathbf{0}$ of~\eqref{eq:face-normal-full}.
Frictional form closure --- i.e.\ $\mathbf{0} \in
\operatorname{int}(\mathcal{W})$ --- is then tested constructively by
searching for a nonzero $\boldsymbol{\lambda}$ that produces zero net wrench
while sitting strictly inside every friction cone, via the linear program
\begin{equation}
\begin{aligned}
    \underset{\boldsymbol{\lambda},d}{\operatorname{maximize}}
    \quad & d \\[2mm]
    \text{subject to}
    \quad & \boldsymbol{G}\boldsymbol{\lambda} = \boldsymbol{0},\\
    & F\boldsymbol{\lambda} - \boldsymbol{1}d \geq \boldsymbol{0},\\
    & d \geq 0,\\
    & \boldsymbol{e}^{T}\boldsymbol{\lambda} \leq n_c.
\end{aligned}
\label{eq:lp2}
\end{equation}

Here, $\boldsymbol{1}$ is a vector of ones with the same number of rows as
$F$. The equilibrium constraint $\boldsymbol{G}\boldsymbol{\lambda}=\boldsymbol{0}$
ensures that the contact forces produce zero resultant wrench. Consequently,
the solution represents an internal force distribution that is in
equilibrium.

The inequality $F\boldsymbol{\lambda} - \boldsymbol{1}d \geq \boldsymbol{0}$
can equivalently be written component-wise as
\begin{equation}
    d \leq (F\boldsymbol{\lambda})_k,
    \qquad
    k=1,\ldots,m,
\end{equation}
where $m$ is the number of rows of $F$. Thus, $d$ is a lower bound on all
contact-constraint margins. Maximizing $d$ therefore maximizes the minimum
margin over all constraints.

A normalization constraint is required because the equilibrium and contact
constraints are homogeneous in $\boldsymbol{\lambda}$. Without such a
constraint, any feasible force distribution could be multiplied by an
arbitrary positive scalar, causing $d$ to increase without bound. The
normalization $\boldsymbol{e}^{T}\boldsymbol{\lambda} \leq n_c$ fixes the
scale of the force distribution. For the ordering in~\eqref{eq:lambda-stacked},
the vector $\boldsymbol{e}$ is constructed such that
\begin{equation}
    \boldsymbol{e}^{T}\boldsymbol{\lambda}
    =
    \sum_{i=1}^{n_c} f_{ni}.
\end{equation}

The optimal value $d^\star$ of~\eqref{eq:lp2} is zero when frictional form
closure does not hold, and strictly positive when it does
($\mathbf{0} \in \operatorname{int}(\mathcal{W})$). How $d^\star$ can be used
as a margin is discussed in \cref{chap:margin}.

% ------------------------------------------------------------------
\subsection{Force Closure}
\label{sec:force-closure-def}

% [Moved from theory.tex, section "Force Closure: Formal Definition"
%  (second paragraph, second half). The intersection symbol was missing
%  in the original.]
Def.~28.10 in \cite{prattichizzo2008grasping} adds
$\mathcal{N}(\boldsymbol{G}) \cap \mathcal{N}(\boldsymbol{J}^\top) = \{\mathbf{0}\}$,
meaning the hand's joints can control every squeeze.

% [Moved from background.tex, subsection "Internal Forces and Force
%  Closure". CHANGED: J -> J^T in the nullspace condition and in LP3, to
%  match Def. 28.10 above (the original used N(J) and J lambda = 0, which
%  is dimensionally inconsistent with lambda). Check against the
%  Handbook.]
The existence of internal forces is not, by itself, a problem. In fact,
internal forces are often necessary to maintain a stable grasp. The relevant
question is whether the internal forces can be appropriately controlled by
the robot. This distinction appears in the force-closure analysis through
the condition
\begin{equation}
    \mathcal{N}(\boldsymbol{G})\cap\mathcal{N}(\boldsymbol{J}^\top)=\{\mathbf{0}\},
    \label{eq:fc-nullspace-condition}
\end{equation}
where $\boldsymbol{J}$ is the hand Jacobian (\cref{sec:hand-jacobian}). The
condition states that there should be no non-zero internal force that
simultaneously belongs to the nullspace of both $\boldsymbol{G}$ and
$\boldsymbol{J}^\top$. Such a force would produce no wrench on the object
while also being uncontrollable through the robot's actuation.

\subsubsection{LP3}
\label{sec:lp3}

This is the purpose of the third step of the approximate force-closure
test. LP3 searches for such a problematic contact-force vector by imposing
\begin{equation}
    \boldsymbol{G}\boldsymbol{\lambda}=\mathbf{0}
    \qquad \text{and} \qquad
    \boldsymbol{J}^\top\boldsymbol{\lambda}=\mathbf{0}.
\end{equation}
If LP3 finds a strictly feasible solution with $d^*>0$, an admissible
non-zero internal force exists in this intersection, and force closure does
not exist. Conversely, if $d^*=0$, no such strictly feasible uncontrollable
internal force is found, satisfying the final condition for force closure
when the preceding tests have also been passed.

\todo[inline]{Write out LP3 in full, in the same form as~\eqref{eq:lp2}
(Handbook Sec.~28.4.2).}
